\chapter{Quantum error correction}\label{ch:qec}

Quantum states are fragile. The coupling to the environment (\cref{ch:open}) and imperfect gates continuously introduce errors, and a quantum computation of any useful length is impossible without correcting them.
With error rates $\sim10^{-3}$ per gate, a circuit of $10^3$ gates already fails most of the time, while Shor's algorithm for RSA-2048 needs $\sim10^{10}$ gates.

\section{Classical redundancy and the three quantum obstacles}
Classical computers protect information by redundancy: a bit is copied several times, $0\to000$, $1\to111$, and a majority vote recovers the correct value if one copy is corrupted. If each copy is flipped independently
with probability $p$, the vote fails only when at least two copies are flipped:
\begin{equation}
  p_{\rm fail}=3p^2(1-p)+p^3=3p^2-2p^3,
\end{equation}
which is smaller than $p$ when $p<1/2$. In the quantum case, this idea seems to face three obstacles:
\begin{enumerate}
  \item the no-cloning theorem forbids copying an unknown state (\cref{sec:nocloning});
  \item measuring the qubits to check them would destroy the superposition;
  \item quantum errors are continuous, since a qubit can be rotated by any small angle, and it seems impossible to correct a continuum of errors with finite resources.
\end{enumerate}
Quantum error correction overcomes all three obstacles \cite{Shor1995}. We illustrate it with the simplest example, the three-qubit bit-flip code.

\section{The three-qubit bit-flip code}
\subsection{Encoding}
The logical qubit is encoded in three physical qubits, $\ket0\to\ket{000}$, $\ket1\to\ket{111}$, so that a general state is encoded as
\begin{equation}
  a\ket0+b\ket1\to\ket\psi=a\ket{000}+b\ket{111}.
\end{equation}
This is not a copy of the state, which would be $(a\ket0+b\ket1)^{\otimes3}$, and it does not contradict the no-cloning theorem. It is an entangled state, prepared by two CNOT gates controlled by the first qubit
and acting on two additional qubits initialized to $\ket0$:
\begin{center}
\begin{quantikz}
\lstick{$a\ket0+b\ket1$} & \ctrl{1} & \ctrl{2} & \qw\\
\lstick{$\ket0$} & \targ{} & \qw & \qw\\
\lstick{$\ket0$} & \qw & \targ{} & \qw
\end{quantikz}
\end{center}

\subsection{Errors}
Suppose that the second qubit undergoes a small unwanted rotation about the $x$ axis, $E_2=\e^{-\ii\epsilon_2\sigma_x^{(2)}}=\cos\epsilon_2-\ii\sin\epsilon_2\,\sigma_x^{(2)}\approx1-\ii\epsilon_2\sigma_x^{(2)}$. Since $\sigma_x$ flips
$\ket0\leftrightarrow\ket1$, to first order in $\epsilon_2$
\begin{equation}
  E_2\ket\psi=\left[a\ket{000}+b\ket{111}\right]-\ii\epsilon_2\left[a\ket{010}+b\ket{101}\right].
\end{equation}
The state is now a superposition of the correct state and of the state with the second qubit flipped.

\subsection{Syndrome extraction}
To find out which qubit was flipped without measuring the qubits themselves, we measure \emph{parities}. Add ancilla qubits in the state $\ket{000}$:
\begin{equation}
  E_2\ket\psi\ket{000}=\left[a\ket{000}+b\ket{111}\right]\ket{000}-\ii\epsilon_2\left[a\ket{010}+b\ket{101}\right]\ket{000}.
\end{equation}
A circuit $S$ of CNOT gates writes the \emph{syndrome} of the error into the ancillas: the parity $q_1\oplus q_2$ of the first two code qubits into the first ancilla, and the parity $q_2\oplus q_3$ into the third ancilla
(two CNOTs each; the second ancilla is not needed and stays in $\ket0$, it is kept here only to follow the original notes). On the code basis states,
\begin{align}
  S\ket{000}\ket{000}&=\ket{000}\ket{000},& S\ket{111}\ket{000}&=\ket{111}\ket{000},\nonumber\\
  S\ket{100}\ket{000}&=\ket{100}\ket{100},& S\ket{011}\ket{000}&=\ket{011}\ket{100},\nonumber\\
  S\ket{010}\ket{000}&=\ket{010}\ket{101},& S\ket{101}\ket{000}&=\ket{101}\ket{101},\nonumber\\
  S\ket{001}\ket{000}&=\ket{001}\ket{001},& S\ket{110}\ket{000}&=\ket{110}\ket{001}.
\end{align}
The syndrome depends only on \emph{which qubit has been flipped}, not on the encoded state: both $\ket{000}$ and $\ket{111}$ give $000$, both $\ket{010}$ and $\ket{101}$ give $101$, and so on. Applied to the erroneous state,
\begin{equation}
  SE_2\ket\psi\ket{000}=\left[a\ket{000}+b\ket{111}\right]\ket{000}-\ii\epsilon_2\left[a\ket{010}+b\ket{101}\right]\ket{101}.
\end{equation}
The error is now entangled with the ancillas, while the coefficients $a$ and $b$ appear identically in both terms.

\subsection{Measurement and correction}
We now measure the ancilla qubits.
\begin{itemize}
  \item With probability $\cos^2\epsilon_2\approx1$, the result is $000$. The state collapses onto $a\ket{000}+b\ket{111}$, and nothing needs to be done.
  \item With probability $\sin^2\epsilon_2\approx\epsilon_2^2$, the result is $101$. The state collapses onto $a\ket{010}+b\ket{101}$, and applying $\sigma_x$ to the second qubit restores $a\ket{000}+b\ket{111}$.
\end{itemize}
In the same way, the syndromes $100$ and $001$ indicate a flip of the first and of the third qubit, which is corrected by applying $\sigma_x$ to that qubit. In all cases the original state is recovered \emph{exactly}.

This simple example shows how the three obstacles are overcome. Redundancy is obtained by entanglement, not by copying. The measurement of the syndrome reveals only which error occurred, never the values of $a$
and $b$, so the superposition is preserved. Finally, the measurement \emph{discretizes} the error: a continuous rotation by an arbitrary small angle is projected either onto ``no error'' or onto a full bit flip,
both of which can be handled. Correcting a discrete set of errors is therefore enough to correct a continuum of them.

\subsection{Stabilizer description}
The parity measurements are measurements of the commuting observables $Z_1Z_2$ and $Z_2Z_3$ (eigenvalue $+1$ for even parity, $-1$ for odd). The code space is the joint $+1$ eigenspace of these two operators,
which are called the \emph{stabilizers} of the code. Each single bit flip $X_j$ anticommutes with a different subset of the stabilizers and therefore produces a distinct pattern of $-1$ outcomes:
$X_1\to(-,+)$, $X_2\to(-,-)$, $X_3\to(+,-)$. Measuring the stabilizers does not disturb the encoded information because the logical operators $\bar X=X_1X_2X_3$ and $\bar Z=Z_1$ commute with them.

\subsection{Limitations}
The three-qubit code corrects a flip of a single qubit. If two qubits are flipped, the syndrome points to the wrong qubit and the correction produces a logical error. For independent errors with probability $p$ per qubit,
the failure probability is reduced from $p$ to $3p^2-2p^3\approx3p^2$, an improvement when $p<1/2$. The code does not protect against phase errors ($\sigma_z$): $Z_j$ commutes with the stabilizers and acts on the code
space as the logical $\bar Z$, flipping the relative sign of $a$ and $b$ undetected.

\section{The phase-flip code}
Phase flips become bit flips in the Hadamard basis $\ket\pm=(\ket0\pm\ket1)/\sqrt2$, since $HZH=X$. The phase-flip code encodes $\ket0\to\ket{{+}{+}{+}}$, $\ket1\to\ket{{-}{-}{-}}$ (the bit-flip encoder followed by $H^{\otimes3}$),
and its stabilizers are $X_1X_2$ and $X_2X_3$. It corrects one phase flip but no bit flips.

\paragraph{Experimental demonstration.} Quantum error correction was first demonstrated in 1998 with nuclear spins in liquid-state NMR \cite{Cory1998}. That experiment implemented the three-qubit code for phase errors
on the carbon and hydrogen nuclei of alanine or trichloroethylene molecules. The fidelity of the stored qubit, as a function of the delay during which it was exposed to natural decoherence, decayed with error correction
as $1-\order(t^2)$ instead of $1-\order(t)$, the signature of first-order protection.

\section{Shor's nine-qubit code}
Concatenating the two codes, Shor obtained the first code able to correct an arbitrary error on any single qubit \cite{Shor1995}:
\begin{gather}
  \ket{0_L}=\frac{(\ket{000}+\ket{111})(\ket{000}+\ket{111})(\ket{000}+\ket{111})}{2\sqrt2},\\
  \ket{1_L}=\frac{(\ket{000}-\ket{111})(\ket{000}-\ket{111})(\ket{000}-\ket{111})}{2\sqrt2}.
\end{gather}
Each block of three is a bit-flip code (stabilizers $Z_1Z_2$, $Z_2Z_3$, $Z_4Z_5$, \dots), protecting against one $X$ error per block. A $Z$ error on any qubit of a block flips the sign of that block; the signs of the three
blocks form a phase-flip code, with stabilizers $X_1X_2X_3X_4X_5X_6$ and $X_4\cdots X_9$.
\begin{theorem}[Discretization of errors]
If a code corrects the errors $X_j$, $Z_j$ and $X_jZ_j\propto Y_j$ on qubit $j$, it corrects \emph{any} error $E_j$ acting on qubit $j$ (any unitary, or any Kraus operator of a channel).
\end{theorem}
\begin{proof}
Expand $E_j=e_0I+e_1X_j+e_2Y_j+e_3Z_j$ on the Pauli basis. The syndrome measurement projects $E_j\ket\psi$ onto one of the four terms, each of which is then corrected; since the recovery is linear, it corrects the sum.
For a channel, apply this to each Kraus operator.
\end{proof}

\section{General theory\starred}
\paragraph{Knill--Laflamme conditions.} A code with projector $P$ onto the code space can correct a set of errors $\{E_a\}$ if and only if
\begin{equation}
  PE_a^\dagger E_bP=c_{ab}P
\end{equation}
for a Hermitian matrix $c_{ab}$ independent of the code state. Intuitively: different errors either map the code space to orthogonal subspaces (they can be distinguished) or act identically on it (they need not be),
and no error reveals information about the encoded state.

\paragraph{Stabilizer codes.} An $[[n,k,d]]$ stabilizer code encodes $k$ logical qubits in $n$ physical qubits as the common $+1$ eigenspace of $n-k$ commuting Pauli operators (the stabilizer generators). It
corrects errors on up to $\lfloor(d-1)/2\rfloor$ qubits, where the \emph{distance} $d$ is the minimal weight of a logical operator. Shor's code is $[[9,1,3]]$, Steane's code (built from the classical Hamming code)
is $[[7,1,3]]$ and allows transversal Clifford gates, and the five-qubit code $[[5,1,3]]$ is the smallest code that corrects an arbitrary single-qubit error: by the quantum Hamming bound, the $1+3n$ possible errors (no error,
or one of 3 Paulis on one of $n$ qubits) need $2(1+3n)\le2^n$ dimensions, which requires $n\ge5$.

\paragraph{Fault tolerance and the threshold theorem.} Error correction is itself performed with noisy gates. Fault-tolerant circuits are designed so that a single fault causes at most one error per code block
(e.g.\ by transversal gates, which act qubit by qubit between blocks). The \emph{threshold theorem} (Aharonov, Ben-Or; Kitaev; Knill, Laflamme, Zurek, 1996--97) states that if the physical error rate $p$ is below a
threshold $p_{\rm th}$, arbitrarily long computations can be performed with an overhead polylogarithmic in the size of the computation. For concatenated codes, each level of encoding maps $p\to cp^2$, so that after $k$ levels
$p_L=p_{\rm th}(p/p_{\rm th})^{2^k}$ with $p_{\rm th}=1/c$: the logical error decreases doubly exponentially.

\paragraph{The surface code.} The most promising code for hardware with nearest-neighbour connectivity is Kitaev's surface code: qubits on a $d\times d$ grid, with 4-body $X$ and $Z$ stabilizers on alternating plaquettes,
measured with ancilla qubits. Its threshold is about $1\%$, and its logical error rate scales as
\begin{equation}
  p_L\approx A\left(\frac p{p_{\rm th}}\right)^{(d+1)/2},
\end{equation}
with about $2d^2$ physical qubits per logical qubit. In 2024, Google demonstrated a surface-code memory below threshold, with the logical error rate decreasing by a factor $\Lambda\approx2$ each time the distance was
increased from $3$ to $5$ to $7$. For $p/p_{\rm th}\approx0.1$ and $p_L\sim10^{-12}$, $d\approx21$--$23$, i.e.\ about $10^3$ physical qubits per logical qubit. Clifford gates are relatively cheap in the surface code; the $T$ gate
requires \emph{magic-state distillation}, which dominates the overhead of fault-tolerant algorithms.

\section{Exercises}
\begin{exercise}
Draw the syndrome-extraction circuit of the bit-flip code with two ancillas. Show that a bit flip on an ancilla during the extraction leads to a wrong correction, and explain why syndromes are measured repeatedly in practice.
\end{exercise}
\begin{exercise}
Show that the error $Z_1Z_2$ is not detected by the Shor code and acts trivially on the code space. (Such errors are harmless: the Shor code is \emph{degenerate}.)
\end{exercise}
\begin{exercise}
With $p=10^{-3}$, $p_{\rm th}=10^{-2}$ and $A=0.1$, what surface-code distance gives $p_L<10^{-12}$? How many physical qubits are needed for $1000$ logical qubits?
(Answer: $(d+1)/2\ge11$, $d=21$; about $2\times21^2\times1000\approx9\times10^5$.)
\end{exercise}
